\par\noindent\textbf{(1)} 均匀球表面壳受到球内全部质量 \(M\) 的Newton引力，故 \(\ddot R=-GM/R^2\)。从初始静止 \(R=R_0\) 的能量式 \(\tfrac12\dot R^2-GM/R=-GM/R_0\)，给坍缩支 \(\dot R=-\sqrt{2GM(1/R-1/R_0)}\)。
\par\noindent\textbf{(2)} 令 \(R=R_0\cos^2\eta\)，\(0\le\eta\le\pi/2\)，由动力学给 \(dt/d\eta=2\sqrt{R_0^3/(2GM)}\cos^2\eta\)。积分可得 \(t(\eta)=\sqrt{R_0^3/(2GM)}(\eta+\tfrac12\sin2\eta)\)。
\par\noindent\textbf{(3)} 令 \(\eta=\pi/2\) 得 \(t_{\rm ff}=\frac\pi2\sqrt{R_0^3/(2GM)}\)。由于初始密度 \(\rho_0=3M/(4\pi R_0^3)\)，得到 \(t_{\rm ff}=\sqrt{3\pi/(32G\rho_0)}\)。时间由初始平均密度控制，与绝对球体大小无关。
\par\noindent\textbf{(4)} 均匀、严格同向坍缩时每个质量壳的 \(M(r_0)\propto r_0^3\)，其自由落体时间相同，因而在理想模型内同时抵达中心。真实气体的压力梯度、角动量、磁场和非均匀密度会破坏同步性，中心奇点不能直接等同于可观测的物理终态。
